\chapter{\MakeUppercase{Project Description and Goals}}
\section{Literature review}
A literature review is a structured and critical examination of existing research, studies, theories, and publications related to a specific topic or problem. In a computer science project report, it summarizes and analyzes what other researchers and developers have already done in the same or closely related area.

It is not just a collection of summaries. Instead, it demonstrates your understanding of the field, identifies patterns or trends, compares different approaches, and highlights gaps that your project aims to address.

A literature review is a text of a scholarly paper, which contains the current knowledge including substantive findings, as well as theoretical and methodological contributions to a particular topic. Literature reviews are secondary sources, and do not report new or 
original experimental work. Most often associated with  academic-oriented literature, such reviews are found in academic journals, and are not to be confused with book reviews that may also appear in the same publication. Literature reviews are a basis for research in nearly every academic field (Waldron, 2008b). 

A narrow-scope literature review may be included as part of a peer-reviewed journal article  presenting new research, serving to situate the current study within the body of the relevant 
literature and to provide context for the reader. In such a case, the review usually precedes the methodology and results sections of the work. Producing a literature review may also be part of graduate and post-graduate student work, including in the preparation of a thesis, 
dissertation, or a journal article. Literature reviews are also common in a research proposal or prospectus (the document that is approved before a student formally begins a dissertation or thesis) (Doan et al., 2002; Duzdevich et al., 2014; Ganesh et al., 2016).

The main types of literature reviews are: evaluative, exploratory, and instrumental (Duzdevich et al., 2014; Neil and David, 2016). 
A fourth type, the systematic review, is often classified separately, but is essentially a literature review focused on a research question, trying to identify, appraise, select and synthesize all high-quality research evidence and arguments relevant to that question. A meta-analysis is typically a systematic review using statistical methods to effectively 
combine the data used on all selected studies to produce a more reliable result \cite{duzdevich2014dna}, \cite{mohanaprasad2017optimized}, \cite{gilbarg1977elliptic}, \cite{haykin2004kalman}, \cite{haykin2005cognitive}, \cite{knuth1986computers}

\section{Research Gap}
One of the most important parts of a report. Explain
\begin{itemize}
	\item What current systems fail to address
	\item Limitations in previous studies
	\item Missing features
	\item Performance weaknesses
	\item Scalability or usability issues
\end{itemize}

\paragraph{}

\section{Objectives}
\paragraph{}

\section{Problem statement}
\paragraph{}

\section{Project plan}
\paragraph{}


